\subsection{Transformación de datos}

Con el conjunto ya seleccionado y limpio, la pregunta que sigue no es qué
transformación aplicar sino si hace falta aplicar alguna. Transformar una
variable tiene un costo: los coeficientes dejan de leerse en las unidades
originales y, en el caso de la discretización, se pierde información de manera
irreversible. Por ello cada decisión de esta sección se sostiene en una medición
hecha sobre nuestros propios datos, y las transformaciones que se descartaron se
documentan junto con su razón.

El diagnóstico se organiza en tres preguntas: si las variables están en escalas
comparables, qué tan lejos se encuentran de la simetría y si sus relaciones son
lineales. Las tres importan porque las técnicas que siguen ---componentes
principales, análisis factorial, conglomerados y discriminante--- operan sobre la
matriz de covarianza o sobre distancias, y solo capturan la componente lineal de
la asociación.

\subsubsection*{Escalas: la varianza está mal repartida}

Las variables medidas conviven en unidades incomparables: ppb, ppm,
$\mu$g/m$^{3}$, grados Celsius, mmHg y km/h. La consecuencia no está en las
medias sino en las varianzas. En escala cruda, \texttt{PM10} y \texttt{NOX}
concentran el $71\%$ de la varianza total del conjunto, mientras que \ce{CO}
aporta apenas el $0.02\%$; no porque sea irrelevante para el ozono, sino porque
se mide en ppm y sus números son pequeños. Un análisis de componentes principales
sobre la matriz de covarianza de estos datos no descubriría estructura alguna:
devolvería una primera componente que es esencialmente \texttt{PM10} y una
segunda que es \texttt{NOX}.

Se estandarizan por tanto todas las variables numéricas mediante puntuación
$z$, es decir, $z=(x-\bar{x})/s$. Conviene subrayar que estandarizar equivale
exactamente a correr el análisis sobre la \emph{matriz de correlación} en lugar
de la de covarianza; es la misma decisión enunciada de dos maneras. Tras la
estandarización, cada una de las once variables del conjunto de trabajo aporta
el $9.09\%$ de la varianza total, que es el reparto uniforme buscado.

Se consideraron dos alternativas y ambas se descartaron. El escalamiento
\textit{min-max} queda determinado por el mínimo y el máximo, los valores más
frágiles de este conjunto: como la limpieza conservó deliberadamente los
extremos, un solo registro atípico comprimiría el resto de la variable en una
franja diminuta del intervalo $[0,1]$. El escalamiento robusto se descarta por
una razón de fondo: los valores extremos aquí no son ruido, son el fenómeno que
se quiere modelar, y restarles peso trabajaría en contra del objetivo de
clasificar episodios de concentración alta.

\subsubsection*{Forma: asimetría y linealidad}

Salvo la temperatura y la presión, todas las variables medidas presentan
asimetría positiva, y no por accidente: una concentración es una cantidad no
negativa con cola larga, pues hay muchas horas de aire limpio y pocas de
episodio, pero esas pocas se alejan mucho. La prueba de Shapiro-Wilk confirma que
ninguna es normal, con estadísticos que van de $W=0.60$ en \texttt{WSR} a
$W=0.99$ en \texttt{PRS}.

Para elegir la transformación se ajustó el parámetro $\lambda$ de Box-Cox por
máxima verosimilitud en cada variable. El resultado, mostrado en el
\cref{tab:transformacion}, es notablemente consistente: todos los valores caen
entre $-0.10$ y $0.35$, es decir, alrededor de cero. Los datos están pidiendo un
logaritmo. Se adopta $\ln(1+x)$ en lugar de la transformación de Box-Cox
propiamente dicha por tres motivos: deja la asimetría igualmente dentro de la
banda de simetría aproximada, admite los ceros legítimos que Box-Cox prohíbe
---\ce{NO2} y \ce{PM2.5} reportan cero en horas limpias, y descartarlos sería
inventar datos faltantes--- y utiliza un mismo exponente para todas las
variables, lo que las mantiene comparables entre sí.

\begin{table*}[t]
  \centering
  \caption{Diagnóstico y efecto de la transformación logarítmica sobre las
    variables asimétricas del conjunto horario. La correlación con el \ce{O3} se
    contrasta con el coeficiente de Spearman, que mide la asociación monótona
    con independencia de la escala y funciona por tanto como techo de lo
    recuperable.}
  \label{tab:transformacion}
  \footnotesize
  \begin{tabular}{lccccccc}
    \toprule
    & \multicolumn{2}{c}{Asimetría} & & \multicolumn{3}{c}{Correlación con \ce{O3}} \\
    \cmidrule(lr){2-3} \cmidrule(lr){5-7}
    Variable & Original & Con $\ln(1+x)$ & $\lambda$ de Box-Cox & Original & En escala logarítmica & Spearman \\
    \midrule
    \ce{O3}    & 1.29 & $-0.51$ & 0.35    & --- & --- & --- \\
    \ce{CO}    & 1.82 & 0.27    & 0.30    & $-0.075$ & $-0.122$ & $-0.107$ \\
    \ce{NO2}   & 1.93 & $-0.18$ & 0.20    & $-0.272$ & $-0.359$ & $-0.363$ \\
    NOX        & 3.63 & 0.45    & $-0.10$ & $-0.329$ & $-0.470$ & $-0.434$ \\
    \ce{PM10}  & 3.66 & $-0.27$ & 0.15    & 0.019 & $-0.040$ & $-0.002$ \\
    \ce{PM2.5} & 5.81 & $-0.12$ & 0.10    & 0.000 & $-0.078$ & $-0.045$ \\
    \texttt{WSR} & 6.94 & $-0.33$ & 0.25  & 0.274 & 0.448 & 0.482 \\
    \bottomrule
  \end{tabular}
\end{table*}

El argumento decisivo, sin embargo, no es la normalidad sino la linealidad. Las
relaciones entre contaminantes son de naturaleza multiplicativa, y el logaritmo
las convierte en aditivas. En el \cref{tab:transformacion} se observa que en
escala original el coeficiente de Pearson se queda muy por debajo del de
Spearman, mientras que en escala logarítmica casi lo alcanza: la correlación
entre \ce{O3} y NOX pasa de $-0.329$ a $-0.470$, y la del \ce{O3} con la
velocidad del viento, de $0.274$ a $0.448$. Esa diferencia no es cosmética, es
señal que el análisis de componentes principales y la regresión no habrían
visto. El signo negativo del NOX frente al ozono es el esperado a escala horaria:
el \ce{NO} titula al ozono en las horas de tráfico.

No se transforman la temperatura ni la presión, cuyas asimetrías de $-0.42$ y
$0.10$ las sitúan ya dentro de la banda de simetría aproximada; transformarlas
solo dificultaría leer los coeficientes, que están en grados Celsius y mmHg.
Tampoco se transforman las componentes vectoriales del viento obtenidas en la
limpieza, porque su signo codifica la dirección y cualquier transformación no
lineal rompería la interpretación vectorial y la simetría alrededor de cero que
las hace utilizables.

\subsubsection*{Discretización}

La discretización se aplicó de forma selectiva, únicamente donde el intervalo
significa algo. La precipitación es el caso claro: con un $98.57\%$ de horas en
cero, no es una variable continua asimétrica sino un proceso de ocurrencia, y el
logaritmo apenas reduce su asimetría de $238$ a $27.7$. Además, el mecanismo que
interesa para el ozono no es cuántos milímetros cayeron sino que hubo nubosidad y
depósito húmedo. Se conserva la variable continua y se agrega una indicadora de
ocurrencia de lluvia.

La hora y el mes plantean un problema distinto. Como enteros, mienten sobre la
distancia entre categorías: colocan las 23:00 a veintitrés unidades de las 00:00
cuando son consecutivas. La solución no es agruparlas en bloques sino
proyectarlas al círculo mediante $\sin(2\pi h/24)$ y $\cos(2\pi h/24)$, y de
manera análoga para el mes, con lo que diciembre y enero recuperan su vecindad.

La variable objetivo, en cambio, no admite discretización libre: sus umbrales los
fija la NOM-020-SSA1-2021. Sobre los $23\,931$ días-estación, el umbral vigente
de $65$ ppb marca como excedencia el $10.65\%$ de los casos; el de $60$ ppb, el
$16.04\%$, y el de $51$ ppb, el $30.70\%$. El desbalance del primer umbral tiene
una implicación práctica para el modelado: un clasificador que respondiera
siempre \enquote{no excede} acertaría el $89\%$ de las veces sin aportar nada, de
modo que la exactitud no puede emplearse como métrica de evaluación y conviene
tomar el umbral de $51$ ppb como escenario principal.

\subsubsection*{Atributos derivados}

El conjunto limpio tiene resolución horaria, pero la variable objetivo, el MDA8,
es diaria. Ambas tablas no se pueden cruzar sin agregar antes los predictores a
día-estación, y la forma de agregar cada uno responde a su mecanismo físico y no
a la conveniencia. De la temperatura se toma el máximo diario, porque la
formación fotoquímica escala con el pico y no con el promedio de la jornada. De
los precursores ---NOX, \ce{NO2} y \ce{CO}--- se toma el promedio de la ventana
de 06:00 a 09:00, correspondiente a la emisión matutina del tráfico, que es la
que alimenta el ozono de la tarde; promediar las veinticuatro horas mezclaría esa
emisión con la titulación nocturna y cancelaría justamente la señal de interés.
Las variables de ventilación se promedian sobre el día, por tratarse de un efecto
acumulado sobre la masa de aire, y la precipitación se acumula. El resultado es la
tabla de trabajo, de $23\,931$ observaciones por $45$ variables.

El MDA8 no se transforma: su asimetría es de $0.64$, prácticamente simétrica, ya
que promediar ocho valores consecutivos hace buena parte del trabajo. Además, sus
coeficientes deben permanecer en ppb, que es la unidad en que está escrita la
norma.

\subsubsection*{Dos precauciones}

Se evaluó centrar cada variable dentro de su estación para eliminar el efecto de
sitio. Un análisis de varianza del MDA8 por estación arroja un efecto
estadísticamente significativo ($F=89.8$), pero que explica solo el $5.3\%$ de la
varianza total. Centrar por estación borraría el contraste espacial a cambio de
muy poco, y ese contraste es precisamente lo que el análisis de conglomerados
debe descubrir. La estación se reserva, por tanto, para validar el agrupamiento
obtenido, no como corrección previa.

Finalmente, los parámetros de estandarización ---centro y escala de cada
variable--- se almacenan por separado del conjunto transformado. Esto responde a
dos necesidades: poder revertir la transformación para reportar los resultados en
unidades originales y, sobre todo, poder recalcularlos únicamente sobre el
subconjunto de entrenamiento cuando se realice la partición para el modelado. De
calcularse sobre la totalidad de los datos, el subconjunto de prueba filtraría
información sobre su propia distribución hacia el modelo.
